\subsection{Consistency of the dominant-component response}\label{app:component-consistency}
The comparison covers 42 selected configurations and 126 checkpoints across all six methods, all tested language capacities and the largest FLUX budget for each method. Restoring dominant components gives the largest mean reduction in general-text NLL or general-image drift in all 42 configurations and a larger reduction than restoring either intermediate or trailing components in 123/126 checkpoints. These counts support the recurring dominant-component forgetting response across models and tasks.

\paragraph{Consistency with smaller edits.}
We measure edit norm as the square root of the summed squared Frobenius changes from trained weights across all adapted matrices, using the stored weights. We select cases where restoring intermediate or trailing components produces a larger edit norm than restoring dominant components. Configuration comparisons use seed means of norms and forgetting reductions. Dominant restoration still gives the largest mean forgetting reduction among all three restorations in 17/17 eligible configurations, comprising 8 Qwen, 6 Llama and 3 FLUX configurations. At the checkpoint level, the count is 48/51, with 22/24 for Qwen, 18/18 for Llama and 8/9 for FLUX. Each configuration or checkpoint is counted once, including when both alternatives have larger norms. Requiring both alternatives to have larger norms gives 4/4 configurations and 10/12 checkpoints. The eligible cases involve LoRA, DoRA and MiLoRA. These counts show that the dominant-component advantage can persist with smaller edits, so edit magnitude alone does not determine the ordering of forgetting reductions.
